\documentclass[tikz,border=8pt]{standalone}
\usepackage[UTF8,fontset=windows]{ctex}
\usepackage{tikz}
\usetikzlibrary{arrows.meta, positioning}

\begin{document}
\begin{tikzpicture}[
    font=\small,
    >=Latex,
    line/.style={-Latex, very thick, draw=gray!70},
    dashline/.style={-Latex, thick, draw=gray!60, dashed},
    box/.style={draw, rounded corners=5pt, very thick, align=center, minimum height=11mm},
    src/.style={box, fill=violet!25, draw=violet!70!black, text width=4.1cm},
    io/.style={box, fill=cyan!18, draw=cyan!55!black, text width=4.2cm},
    proc/.style={box, fill=green!22, draw=green!45!black, text width=4.4cm},
    outbox/.style={box, fill=gray!20, draw=gray!65!black, text width=4.5cm},
    note/.style={draw=none, align=left, font=\scriptsize, text width=5.4cm}
]

\node[font=\Large\bfseries] (title) at (0, 1.2) {LTE 下行信道完整译码流程};
\node[src] (sync) at (0, 0) {小区搜索与同步\\PSS/SSS、PCI、帧定时};

% PBCH lane
\node[io]   (pbchre) at (-8.2, -2.4) {PBCH 资源提取与均衡\\中心 72 子载波};
\node[proc] (pbchllr) at (-8.2, -4.2) {QPSK 软比特 / LLR};
\node[proc] (pbchdec) at (-8.2, -6.4) {PBCH 译码\\解扰 $\rightarrow$ 速率恢复\\尾咬卷积译码 $\rightarrow$ CRC};
\node[outbox]  (mib) at (-8.2, -8.8) {MIB 输出\\带宽 / PHICH / SFN 高位};

% PDCCH lane
\node[io]   (cfi) at (0, -2.4) {控制区大小\\PCFICH/CFI 或外部可信输入};
\node[io]   (pdcchre) at (0, -4.5) {PDCCH 有效 RE 提取与均衡\\排除 CRS / PCFICH / PHICH};
\node[proc] (td) at (0, -6.6) {2Tx 发射分集去映射\\1Tx 时可直接跳过};
\node[proc] (pdcchsoft) at (0, -8.5) {QPSK 软比特};
\node[proc] (search) at (0, -10.7) {REG / CCE 重组 + 盲检索\\候选聚合级别 / 搜索空间};
\node[proc] (dcidec) at (0, -13.0) {候选译码\\解扰 $\rightarrow$ 速率恢复\\卷积译码 $\rightarrow$ CRC-RNTI};
\node[outbox]  (dci) at (0, -15.5) {DCI 输出\\RNTI / 格式 / CCE / 聚合级别 / 调度授权};

% PDSCH lane
\node[io]   (grant) at (8.2, -4.2) {PDSCH 调度构建\\MIB + DCI + 小区上下文};
\node[io]   (pdschre) at (8.2, -6.6) {PDSCH RE 提取与均衡\\PRB / 起始符号 / 参考信号排除};
\node[proc] (layer) at (8.2, -8.9) {层处理 / 传输模式恢复\\层映射 / 分集 / 码字路径};
\node[proc] (pdschsoft) at (8.2, -11.1) {星座软比特};
\node[proc] (pdschdec) at (8.2, -13.6) {PDSCH 译码\\解扰 $\rightarrow$ 速率恢复\\Turbo 译码 $\rightarrow$ CRC};
\node[outbox]  (tb) at (8.2, -16.0) {传输块输出\\TB0 / TB1};
\node[outbox]  (mac) at (8.2, -18.2) {MAC PDU 输出\\子头 / SDU / 控制单元};
\node[outbox]  (higher) at (8.2, -20.4) {更高层输出\\RRC / NAS / IP 载荷};

\draw[line] (sync) -- (pbchre);
\draw[line] (pbchre) -- (pbchllr);
\draw[line] (pbchllr) -- (pbchdec);
\draw[line] (pbchdec) -- (mib);

\draw[line] (sync) -- (cfi);
\draw[line] (cfi) -- (pdcchre);
\draw[line] (pdcchre) -- (td);
\draw[line] (td) -- (pdcchsoft);
\draw[line] (pdcchsoft) -- (search);
\draw[line] (search) -- (dcidec);
\draw[line] (dcidec) -- (dci);

\draw[line] (mib.east) -- ++(1.8,0) |- (grant.west);
\draw[line] (dci.east) -- ++(1.5,0) |- (grant.west);
\draw[line] (grant) -- (pdschre);
\draw[line] (pdschre) -- (layer);
\draw[line] (layer) -- (pdschsoft);
\draw[line] (pdschsoft) -- (pdschdec);
\draw[line] (pdschdec) -- (tb);
\draw[line] (tb) -- (mac);
\draw[line] (mac) -- (higher);

\node[note] at (0, -23.2) {工程实现边界说明：完整 LTE 接收机应产出 MIB、DCI、TB 与 MAC PDU；\\当前 \texttt{MMSE\_CPP} 主要覆盖资源提取、CRS 信道估计、MMSE 均衡，以及 PDCCH 的 2Tx 发射分集去映射。};

\end{tikzpicture}
\end{document}
